%%%%%%%%%%  scan idx 359  |  folio 346  %%%%%%%%%%
On le
munit d'une structure de \(K\)-\(\lambda\)-algèbre de la façon suivante : nous avons remarqué plus haut que l'image du foncteur \(\gamma'\colon\widehat G^t\longrightarrow\widehat G^u\circ i\circ\widehat G^t_g\) est dans le sous-groupe des séries formelles en \(u\) dont les coefficients en degré non nul sont dans \(\widehat G^t\). Utilisant le morphisme canonique : \(\widehat G\longrightarrow 1+\widehat{\operatorname{Gr}\,\cdot\,}^{+}\), qui donne : \(\widehat G\longrightarrow K\times(1+\widehat{\operatorname{Gr}\,\cdot\,}^{+})\), on en déduit finalement un morphisme :
\[
  \widehat G^t\xrightarrow{\gamma'}\widehat G^u\circ i\circ(K\times(1+\widehat{\operatorname{Gr}\,\cdot\,}^{+})),
\]
qui d'après 1.12 définit un morphisme de foncteurs sur\qquad:
\[
  1+\widehat{\cdot}{}^{+}\xrightarrow{\gamma'}\widehat G^u\circ i\circ(K\times(1+\widehat{\cdot}{}^{+})).
\]
Les polynômes universels définissant ce morphisme sont les coefficients des \({\gamma'}^{i}(1+X_1+\dots+X_n+\dots)\)~; le \(j\)-ième coefficient est la composante isobare de poids \(j\) de \(Q'_{i,j}\) (démonstration analogue à celle de 1.15), et sera notée \(Q''_{i,j}\).

On étend enfin le morphisme ainsi construit à \(K\times(1+\widehat{\cdot}{}^{+})\) de façon naturelle, en utilisant la \(\lambda\)-structure de \(K\). On obtient donc une structure de pré-\(\lambda\)-anneau sur \(K\times(1+\widehat{\cdot}{}^{+})\), définie par les opérations \(\gamma\) correspondantes, qui sont telles que :
\begin{equation}\tag{6.1.4}
  \gamma_t(a,\varphi)=\gamma_t(a).\Bigl(1+\sum_{p\ge1}{\gamma'}^{p}(\varphi)t^{p}\Bigr),
\end{equation}
et qui seront notées \(\gamma\) au lieu de \(\gamma'\), aucune confusion n'étant possible.

Pour vérifier que c'est une structure de \(\lambda\)-anneau, on vérifie la commutativité de (3.7.2)~; elle est connue pour les éléments de \(K\)~; pour les éléments de \(1+\widehat{\cdot}{}^{+}\), elle résulte de celle de (3.7.2) pour \(\widehat G\) par passage au quotient. On en tire aussitôt le résultat.



% Local macros for Exposé V, §§6--7
\providecommand{\Ghat}{\widehat G}
\providecommand{\Aqhat}{\widehat A_{\mathbb Q}}
\providecommand{\Ch}{\operatorname{Ch}}
\providecommand{\Log}{\operatorname{Log}}
\renewcommand{\id}{\operatorname{id}}
\providecommand{\gr}{\operatorname{gr}}
\providecommand{\Hom}{\operatorname{Hom}}
\providecommand{\epsi}{\varepsilon}
\providecommand{\clgr}{\operatorname{cl}}
\providecommand{\Fil}{\operatorname{Fil}}
\providecommand{\Gr}{\operatorname{Gr}}
\providecommand{\Sym}{\operatorname{Sym}}


